\documentclass[10pt,a4paper]{amsart}
\usepackage[margin=19mm]{geometry}
\usepackage{amsmath,amssymb,mathtools}
\usepackage[scr=rsfs,cal=euler]{mathalfa}
\usepackage{xfrac}
\usepackage[unicode,hidelinks,bookmarks=false]{hyperref}
\newcommand{\ie}{i.e.\ }
\newcommand{\cf}{cf.\ }
\newcommand{\resp}{respectively\ }
\newcommand{\Subsection}{\subsection}
\providecommand{\nobreakdash}{\nobreak\hbox{-}}
\setlength{\emergencystretch}{2em}
\multlinegap0pt
\usepackage{amssymb}
\usepackage{enumerate}
\usepackage{mathtools}
\usepackage[shortlabels]{enumitem}
\newtheorem{lemma}{Lemma}
\title{An Elementary Analysis of the Prime Partition Function}
\author{Asaf Cohen Antonir, Asaf Shapira}
\date{Source: arXiv:2606.06068v1 (CC BY 4.0)}
\begin{document}
\maketitle

\noindent\emph{Source and licence.} An Elementary Analysis of the Prime Partition Function. Version arXiv:2606.06068v1 (CC BY 4.0), distributed by arXiv under the Creative Commons Attribution 4.0 International licence, http://creativecommons.org/licenses/by/4.0/, as recorded in the arXiv licence metadata for this exact version and shown on the abstract page https://arxiv.org/abs/2606.06068v1. Full licence notice: \texttt{licenses/arxiv-2606.06068-CC-BY-4.0.txt}.\par\smallskip

\noindent\textit{Editorial scope.} Counted unit: the article's lemma lem:prime\_sum (source0.tex lines 394--396). The article's complete original proof of it is delivered with it (source0.tex lines 551--598, one \texttt{proof} environment of 48 lines, which the article places away from the statement). No mathematics is written, repaired or re-derived: the statement and the proof are the article's own lines, in the article's own notation, and no retained line is substituted or re-flowed.  No further source-local context is needed: every cross-reference of every retained line resolves inside the two delivered spans.  Nothing was omitted from any retained span, so the package carries no omission record.  No retained line contains a citation, so the wrapper carries no bibliography and no bibliography entry.  No retained line contains a figure environment, an image-inclusion command or drawing code, so no image is delivered and the package declares no asset dependency.  Editorial formatting only, in addition: an AMS \texttt{amsart} preamble that restates the article's own declarations and macros used by the retained lines, namely the environments \texttt{lemma} on separate counters, together with the standard packages those lines need.  The retained environments print in this document as Lemma 1 in order of appearance, whereas the article prints the same items with the numbering of its own full document.  The complete native source of the article as distributed in the arXiv e-print for this exact version is bundled unchanged as \texttt{sources/202155/source0.tex}.
\begin{lemma}\label{lem:prime_sum}
We have $\sum_{p\in \mathcal{P}} p e^{-\delta p} \sim {\delta^{-2}}/{\log(1/\delta)}$ as $\delta \to 0$.
\end{lemma}

\begin{proof}[Proof of Lemma \ref{lem:prime_sum}]
We start by showing the upper bound, that is, for every $\varepsilon>0$ there exists $\delta_0>0$ such that for all $\delta<\delta_0$ we have
\[
    \sum_{p\in\mathcal{P}}pe^{-p\delta}\leq \frac{1+\varepsilon}{\delta^2\log(1/\delta)}\;.
\]

First, we show that for every $\xi>0$ there exists $\eta_0>0$ such that for all $\eta<\eta_0$ the following holds:
\begin{equation}\label{eq:evaluating_the_integral_primes_upper}
        \int_{1}^{\infty}x\log(x)e^{-\eta x\log(x)}dx\leq \frac{1+\xi}{\eta^2\log(1/\eta)}\;.
\end{equation}
Indeed, let $0<\xi<1$ and let $a>0$ be sufficiently large so that $\log(x)+1\geq (1-\xi/8)\log(x\log(x))$ for all $x\geq a$. Then,
\begin{align*}\label{eq:evaluation_with_integral}
    \int_{1}^{\infty}x\log(x)e^{-\eta x\log(x)}dx &\leq \int_{1}^{a}x\log(x)e^{-\eta x\log(x)}dx+\frac{1}{(1-\xi/8)\eta^2}\int_{\eta a\log(a)}^{\infty}\frac{z e^{-z}}{\log(z/\eta)}dz\nonumber\\
    &\leq \int_{1}^{a}x^2dx+\frac{1+\xi/4}{\eta^2}\left(\int_{\eta a\log(a)}^{1/\log^2(1/\eta)}\frac{ze^{-z}}{\log(z/\eta)}dz+\int_{1/\log^2(1/\eta)}^{\infty}\frac{ze^{-z}}{\log(z/\eta)}dz\right)\nonumber\\
    &\leq \frac{\xi/2}{\eta^2\log(1/\eta)}+\frac{1+\xi/4}{\eta^2}\left(\frac{C}{\log^2(1/\eta)}+\frac{1}{\log(1/\eta)-2\log\log(1/\eta)}\int_{0}^{\infty}ze^{-z}dz\right)\\
    &\leq \frac{1+\xi}{\eta^2\log(1/\eta)}\;,
\end{align*}
where the first inequality holds by the substitution $z=\eta x \log x$, which, by the choice of $a$, guarantees that
$dz=\eta(\log(x)+1)dx\geq \eta(1-\xi/8)\log(z/\eta)dx$, the third inequality holds for some constant $C$ that depends only on $a$, and provided $\eta_0$ is small enough, the last inequality holds provided $\eta$ is sufficiently small and the fact that $\int_{0}^{\infty}ze^{-z}dz=1$.

Next, for every $\delta,x>0$ let $f_\delta(x)=x\log(x)e^{-\delta x\log(x)}$. We now show that for any $\xi>0$ there exists $\eta_0>0$ so that for all $\eta<\eta_0$ we have
\begin{equation}\label{eq:sum_to_int}
    \sum_{m=2}^{n} f_\eta(m)\leq \int_{1}^{\infty}f_\eta(x) dx+\frac{\xi}{\eta^2\log(1/\eta)}\;.
\end{equation}
Choosing the parameters appropriately, the prime number theorem, \eqref{eq:evaluating_the_integral_primes_upper}, and \eqref{eq:sum_to_int} assert that
\[
    \sum_{p\in \mathcal{P}}pe^{-\delta p}\sim \sum_{m=2}^{\infty}f_\delta(m)\leq \int_{1}^{\infty}f_{\delta}(x)dx+\frac{\varepsilon/2}{\delta^2\log(1/\delta)}\leq \frac{1+\varepsilon}{\delta^2\log(1/\delta)} \;.
\]
which finishes the proof of the upper bound. To prove \eqref{eq:sum_to_int}, note that for any fixed $\eta$ we have
$\frac{d}{dx}f_\eta(x)=-(\log(x)+1)(\eta x\log(x)-1)e^{-\eta x\log(x)}$.
Hence, provided $\eta$ is small enough and letting $b\sim \eta^{-1}/\log(1/\eta)>1/e$ be such that $f_\eta(x)$ is maximal, we have the following as $f_\eta(x)$ is increasing from $1/e$ to $b$, and decreasing from $b$ to infinity:
\[
    \sum_{m=2}^{\lfloor b \rfloor -1}f_\eta(m)\leq \int_{1}^{\lfloor b\rfloor}f_\eta(x)dx\quad \text{and}\quad \sum_{m=\lfloor b \rfloor+2}^{\infty} f_\eta(m)\leq \int_{\lfloor b\rfloor}^{\infty}f_\eta(x) dx\;.
\]

We see that the integral $\int_{0}^{\infty}f_\eta(x)dx$ accounts for $\sum_{m=1}^{\infty}f_\eta(m)$ except possibly for $f(\lfloor b\rfloor)$ and $f(\lfloor b\rfloor +1)$, hence, we obtain \eqref{eq:sum_to_int} as follows
\begin{align*}
    \sum_{m=2}^{\infty}f_\eta(m)&\leq \int_{1}^{\infty}f_\eta(x)dx+2f_\eta(b)=\int_{1}^{\infty}f_\eta(x)dx+\frac{2}{\eta e}\leq \int_{1}^{\infty}f_\eta(x)dx+\frac{\xi}{\eta^2\log(1/\eta)}\;,
\end{align*}
where the first inequality holds by the above inequalities and the definition of $b$ being the maximal point of $f_\eta(x)$, and the second inequality holds provided $\eta$ is sufficiently small.

The proof of the lower bound is almost identical, therefore, we will only sketch it. One can obtain analogous inequalities to \eqref{eq:evaluating_the_integral_primes_upper} and \eqref{eq:sum_to_int}, that is
\[
    \int_1^{\infty}f_\eta(x)dx\geq \frac{1-\xi}{\eta^2\log(1/\eta)}\quad \text{and}\quad
    \sum_{m=2}^{n} f_\eta(m)\geq \int_{1}^{\infty}f_\eta(x)dx-\frac{\xi}{\eta^2\log(1/\eta)}\;.
\]
Using the prime number theorem and the above inequalities, we derive the lower bound, similar to the upper bound.
\end{proof}

\end{document}
